\section{Introduzione al capitolo}\label{sec:intro_cap2}

Il capitolo esplora il background scientifico e tecnologico alla base del lavoro proposto.
Partiamo dal \textit{problema fondamentale}: come fornire accesso content-aware a repository RDF di dimensioni massive per utenti privi di competenze tecniche formali?

La risposta richiede una comprensione su tre livelli di complessità:
\begin{enumerate}
    \item \textbf{Lo stato dell'arte}: come gli approcci precedenti (form rigidi, endpoint SPARQL nudi) falliscono nel dominio umanistico;
    \item \textbf{I fondamenti teorici}: i principi dell'architettura dell'informazione che guidano un'alternativa;
    \item \textbf{Le tecnologie emergenti}: come Retrieval-Augmented Generation (RAG) e Large Language Models (LLM) abilitano una soluzione più flessibile.
\end{enumerate}

Procederemo così: nel \cref{sec:statoarte} analizziamo il divario tra i bisogni esplorativi degli studiosi umanisti e la rigidità sintattica dei sistemi tradizionali. Successivamente, nel \cref{sec:archinfo}, esaminiamo i principi fondativi dell'architettura dell'informazione che giustificano il ricorso agli LLM.\@
Infine, nel \cref{sec:rag}, descriviamo il paradigma RAG come soluzione tecnologica che bilancia flessibilità e affidabilità.

\section{Stato dell'arte: dall'interrogazione formale alla ricerca semantica}\label{sec:statoarte}
\subsection{Il contesto: Digital Humanities e la barriera della ricerca}\label{subsec:dh_barrier}

Nel dominio delle Digital Humanities (DH) e della gestione del patrimonio culturale, la ricerca all'interno di repository RDF rappresenta un ostacolo significativo per chi manca di competenze tecniche avanzate.
Come evidenziato in letteratura, lavorare con strumenti come SPARQL impone vincoli rigorosi\cite{Vas25} e richiede una familiarità profonda con ontologie complesse e convenzioni di metadati\cite{Vas24}. 

Tali approcci si basano su query create manualmente, che risultano fragili e difficili da adattare a esigenze mutevoli\cite{Mar25}. Questo pone il ricercatore di fronte all'ostacolo cognitivo della ``scatola vuota'': chi possiede un bisogno informativo vago o in via di definizione viene paralizzato dalla necessità di dichiarare tutto a priori.

\subsection{Il divario tra ricerca strutturata e ricerca esplorativa}\label{subsec:berrypicking}

Alla base dei tradizionali sistemi di Information Retrieval (IR) vi è un assunto fondamentale: che l'utente inizi l'interazione con un bisogno informativo perfettamente strutturato.
Nel dominio umanistico, tuttavia, l'utente procede per esplorazione.

Questo comportamento è stato teorizzato da Marcia Bates con il modello \textit{berrypicking}\cite{Bat89}: una strategia di ricerca non lineare in cui gli studiosi si muovono iterativamente, aggiustando la rotta e raffinando gli obiettivi man mano che acquisiscono nuove informazioni durante l'esplorazione.
In un tale scenario di ricerca esplorativa, assume un valore inestimabile il fenomeno della \textit{serendipity}, la scoperta inattesa di pattern rilevanti durante la navigazione dei dati, che gli strumenti tradizionali non incoraggiano, anzi, spesso ostacolano.

Considerando il netto contrasto tra la rigidità sintattica di SPARQL e la natura fluida ed evolutiva della ricerca umanistica, emerge con chiarezza la necessità di un nuovo paradigma di interazione, che consenta ai ricercatori di esprimere bisogni emergenti senza dover ricorrere a linguaggi formali.

\subsection{La ``scatola vuota'' e le tre barriere di SPARQL}\label{subsec:tre_barriere}

Costringere ricercatori umanisti a esprimere bisogni complessi tramite endpoint SPARQL genera un ostacolo su tre livelli:

\begin{enumerate}
    \item \textbf{Barriera cognitiva}: L'interazione esige la scrittura in un linguaggio dichiarativo formale e l'obbligo di conoscere complessi prefissi ontologici (es.\ \code{dc:}, \code{exif:}, \code{ebucore:}) solo per abbozzare una query di base.
    \item \textbf{Barriera strutturale}: L'assenza di una mappa visiva delle relazioni all'interno del grafo RDF fa sì che l'utente non sappia quali proprietà esistano realmente nel database né come i dati siano formalizzati nell'ontologia.
    \item \textbf{Barriera pratica}: La forte rigidità sintattica determina che un singolo errore di battitura produca un set di risultati vuoto. L'assenza di messaggi di errore parlanti impedisce qualsiasi tentativo di debugging autonomo.
\end{enumerate}

Studi sull'usabilità delle interfacce per Knowledge Graphs confermano che linguaggi come SPARQL impongono vincoli rigidi e richiedono una conoscenza esplicita degli schemi, competenze estranee ai professionisti del patrimonio culturale\cite{Vas25}.
\newpage
\section{Fondamenti di architettura dell'informazione}\label{sec:archinfo}
La progettazione di un sistema di accesso per un database RDF non può limitarsi alla pura efficienza computazionale del recupero dati. Deve invece rispondere ai principi fondamentali dell'\textit{architettura dell'informazione}, ovvero la disciplina che studia come rendere le informazioni realmente trovabili e usabili.

In un contesto umanistico, costringere l'utente a interfacciarsi con la complessità dello schema sottostante attraverso filtri rigidi o linguaggi formali genera un ostacolo cognitivo.
Procediamo ora ad analizzare i principi teorici che giustificano il ricorso a un'interfaccia conversazionale mediata da LLM.\@

    \subsection{L'illusione della ricerca algoritmica}\label{subsec:illusione}

    I sistemi di reperimento basati su query formali (SQL, SPARQL) assumono erroneamente che l'utente sappia sempre cosa cercare e come tradurre il proprio bisogno in un costrutto logico. Peter Morville e Louis Rosenfeld hanno smontato questo assunto come un grave errore concettuale:

    \begin{quotation}
    ``Consideriamo questo modello pericoloso perché basato su un equivoco: che trovare informazioni sia un problema diretto, che può essere risolto da un semplice approccio algoritmico. Dopo tutto, abbiamo risolto la sfida del recupero dati --- che, ovviamente sono fatti e cifre --- con tecnologie database come SQL.\@ Così procede la teoria, si trattano idee e concetti incorporati nei documenti testuali semistrutturati nello stesso modo''\hfill\cite{Mor08}
    \end{quotation}

    I sistemi algoritmici tradizionali falliscono nel dominio umanistico perché sono progettati per recuperare \textit{dati esatti}, mentre gli studiosi ricercano \textit{concetti}, \textit{idee} e \textit{relazioni}. Tali concetti sono per loro natura ambigui e sfumati. Un intermediario LLM colma esattamente questo divario: funge da traduttore semantico che accetta l'ambiguità e la ricchezza del linguaggio naturale e la converte in modo trasparente nell'approccio algoritmico rigido richiesto dalla macchina.

    \subsection{Recall, Precision e il compromesso ineluttabile}\label{subsec:recall_precision}

    Un problema critico delle interfacce tradizionali a filtri è l'impossibilità di calibrare dinamicamente il raggio della ricerca. Morville e Rosenfeld spiegano che ogni ricerca è un compromesso matematico tra due metriche inversamente correlate:

    \begin{quotation}
        \begin{align}
            \var{recall} &= \frac{|\var{doc rilevanti recuperati}|}{|\var{doc nella raccolta}|} \label{eq:recall}\\[5pt]
            \var{precision} &= \frac{|\var{doc rilevanti recuperati}|}{|\var{doc rilevanti della raccolta}|} \label{eq:precision}
        \end{align}
        ``[\ldots] Non sarebbe bello poter ottenere contemporaneamente elevati valori di recall e precision? Purtroppo non si può avere botte piena e moglie ubriaca: recall e precision sono inversamente correlati''\hfill\cite{Mor08}
    \end{quotation}

    In un'interfaccia a maschera con decine di campi e filtri ontologici, l'utente è costretto a dichiarare a priori il livello di precisione desiderato. Se sbaglia un filtro, ottiene zero risultati (eccesso di Precision); se ne omette uno, viene sommerso da migliaia di risultati irrilevanti (eccesso di Recall).

    In una conversazione in linguaggio naturale con un LLM, invece, l'aggiustamento avviene in modo iterativo e fluido. L'utente può iniziare con una richiesta esplorativa e ampia (es.\ ``Mostrami tutti i manoscritti del XIV secolo'', privilegiando il Recall) e, osservando i primi risultati, raffinare il tiro conversando (``Filtra solo quelli con miniature e digitalizzati ad alta risoluzione'', aumentando la Precision). L'LLM permette di variare queste metriche dinamicamente senza costringere l'utente a decostruire e riscrivere manualmente una complessa sintassi di query.

    \subsection{Carico cognitivo: la Legge del minimo sforzo e la Legge di Hick}\label{subsec:cognitive_load}

    Dal punto di vista del carico cognitivo, esporre all'utente l'immensa topologia di un Knowledge Graph attraverso decine di menu a tendina o l'obbligo di conoscere specifiche ontologie porta alla frustrazione potrebbe sfociare nell'abbandono del sistema.

    Luca Rosati analizza questo fenomeno attraverso due principi cognitivi fondamentali. Il primo è la \textit{legge del minimo sforzo}:

    \begin{quotation}
    ``La legge del minimo sforzo (least effort) regola quindi il nostro sistema di information seeking, al punto che spesso accettiamo anche contenuti di bassa qualità o affidabilità (se più facili da ottenere e usare) pur di non ricorrere a strategie di ricerca attive, le quali comportano necessariamente sforzi, tempi e competenze maggiori''\hfill\cite{Ros07}
    \end{quotation}

    Se l'interfaccia richiede la fatica di apprendere prefissi RDF o di comprendere uno schema dati complesso, è plausibile che l'utente rinuncerà alla ricerca.

    Il secondo principio è la \textit{Legge di Hick}, la quale afferma che, in presenza di $n$ alternative equiprobabili, il tempo necessario per selezionarne una è proporzionale (a meno di una costante additiva) al logaritmo in base 2 del numero di scelte più 1:

    \[ \text{tempo} = a + b \log_2(n+1) \]

    I coefficienti $a$ e $b$ dipendono da specifiche condizioni al contorno, tra cui le modalità di presentazione delle opzioni e il grado di familiarità dell'utente~\cite{Ros07}.

    Tuttavia, Rosati sottolinea che, qualora le alternative ``siano presentate in modo confuso'' o costituiscano ``una lista disordinata'' priva di un criterio logico riconoscibile, il tempo di decisione scala in modo lineare, aumentando il carico cognitivo fino a divenire insostenibile e causando paralisi decisionale. Esporre l'utente a un form SPARQL nudo o a una maschera di ricerca avanzata con decine di menu a tendina equivale esattamente a questo: innescare un ``paradosso della scelta'' paralizzante.

    Un'interfaccia conversazionale basata su LLM nasconde questa complessità strutturale nel backend. L'utente si trova di fronte a un singolo ``spazio'' (un box di chat), azzerando l'effetto della Legge di Hick e assecondando la Legge del minimo sforzo. Non dovendo scegliere tra decine di opzioni incomprensibili, il ricercatore delega il carico cognitivo della traduzione ontologica all'LLM, dedicando le proprie energie intellettuali esclusivamente all'analisi e alla scoperta delle risorse.

    \subsection{Paradigma classico: navigazione e filtri rigidi}\label{subsec:paradigma_classico}

    L'architettura dell'informazione si definisce come il design strutturale di ambienti informativi finalizzato a plasmare l'usabilità e la trovabilità\cite{Mor08}. Poggia su quattro sistemi interconnessi:

    \begin{enumerate}[label={(\roman*)}]
        \item sistemi di organizzazione, che categorizzano le informazioni;
        \item sistemi di etichettatura, che ne determinano la rappresentazione linguistica;
        \item sistemi di navigazione, che consentono l'orientamento nello spazio;
        \item sistemi di ricerca, che permettono interrogazioni dirette.
    \end{enumerate}

    Applicato al Datavault, l'interrogazione tradizionale mediante filtri rigidi o SPARQL presenta tre limiti strutturali:

    \begin{enumerate}
        \item Pretende conoscenza a priori di stringhe esatte, incompatibile con la ricerca umanistica\hfill\cite{Vas24};
        \item Non gestisce sinonimia, sfumature e vaghezza del linguaggio naturale;
        \item Obbliga la formulazione di query multi-criterio in sintassi rigide.
    \end{enumerate}

    Queste rigidità abbattono la ``trovabilità'': l'informazione esiste fisicamente ma rimane irraggiungibile per l'utente.

    \subsection{Il nuovo paradigma: ricerca semantica con LLM}\label{subsec:paradigma_semantico}

    L'introduzione di un LLM modifica radicalmente il paradigma di esplorazione dei dati semantici, agendo come un intermediario e traduttore capace di convertire fluidamente il linguaggio naturale in interrogazioni strutturate come query SPARQL.\@

    Questa trasformazione definisce un concetto chiave nell'interazione uomo-dato: non è più il ricercatore a dover imparare il linguaggio della macchina, ma è il sistema stesso che si adatta all'utente, comprendendone le richieste nel formato linguistico naturale. Ne consegue che questo approccio sposta l'intero carico di complessità cognitiva e sintattica dall'utente al sistema computazionale.

    Questo trasferimento della complessità si lega esplicitamente al concetto di ``findability'' (trovabilità)\cite{Ros07}. Lo scopo del nuovo sistema non è semplicemente rendere una risorsa formalmente accessibile, ma garantire che essa risulti effettivamente trovabile dal ricercatore nel momento del bisogno.

    La ricerca semantica mediata dall'LLM aumenta in modo decisivo la trovabilità per tre ragioni:

    \begin{itemize}
        \item abbassa drasticamente la soglia di accesso all'informazione, permettendo a chiunque di fare una domanda diretta al database;
        \item il modello gestisce in modo nativo la vaghezza, l'ambiguità e i sinonimi intrinseci nel linguaggio umano;
        \item consente di esprimere bisogni informativi multi-criterio complessi, impossibili da supportare attraverso form a tendina.
    \end{itemize}

    Mantenendo un doveroso grado di obiettività scientifica, occorre tuttavia dichiarare che il paradigma basato su LLM presenta sfide intrinseche: il rischio di allucinazioni (invenzione di predicati inesistenti), la dipendenza dall'accuratezza del prompt di sistema, e l'elevato costo computazionale delle chiamate ai modelli generativi. Queste metriche verranno analizzate in dettaglio nel \cref{chap:valutazione}.


\section{Retrieval-Augmented Generation (RAG) come soluzione tecnologica}\label{sec:rag}
    \subsection{Il paradigma RAG tradizionale}\label{subsec:rag_trad}

    Retrieval-Augmented Generation (RAG) è un paradigma che combina la potenza dei modelli generativi con l'affidabilità del recupero di informazioni da fonti strutturate. La struttura di base è semplice ma efficace:

    \begin{enumerate}
        \item \textbf{Ingresso}: l'utente pone una domanda in linguaggio naturale;
        \item \textbf{Retrieval}: un modulo di ricerca estrae documenti o frammenti rilevanti da una base di conoscenza strutturata;
        \item \textbf{Augmentation}: il contesto recuperato viene fornito al modello generativo come ``prompt esteso'';
        \item \textbf{Generation}: il modello generativo produce una risposta grounded nei documenti recuperati, riducendo le allucinazioni.
    \end{enumerate}

    Questo paradigma riduce significativamente il rischio di allucinazioni perché il modello è vincolato a operare all'interno di uno spazio di conoscenza ben definito.

    \subsection{RAG applicato a Knowledge Graphs}\label{subsec:rag_kg}

    Nel nostro caso, il ``corpus'' da cui recuperare informazioni è uno \textbf{schema RDF strutturato}. L'applicazione di RAG a Knowledge Graphs richiede una variante: anziché recuperare documenti di testo, recuperiamo \textit{frammenti di schema}, \textit{esempi di predicati}, e \textit{descrizioni di ontologie}.

    Quindi il flusso diventa:

    \begin{enumerate}
        \item \textbf{Ingresso}: domanda dell'utente (es.\ ``Mostrami tutte le immagini digitalizzate con una Nikon nel 2022'');
        \item \textbf{Retrieval}: il sistema accede allo schema RDF del Datavault e recupera i predicati rilevanti (es.\ \code{exif:make}, \code{exif:dateTime});
        \item \textbf{Augmentation}: lo schema e gli esempi di predicati vengono forniti al modello LLM come parte del prompt di sistema;
        \item \textbf{Generation}: il modello LLM produce una query SPARQL coerente con lo schema;
        \item \textbf{Esecuzione}: la query viene eseguita su Blazegraph e i risultati vengono restituiti all'utente.
    \end{enumerate}

    Questo approccio combina i lati positivi di entrambi: la flessibilità del linguaggio naturale e l'affidabilità di una fonte di conoscenza strutturata e verificabile.

    \subsection{Fine-tuning e LoRA per Domain Adaptation}\label{subsec:finetune_lora}

    Un miglioramento ulteriore consiste nel fine-tuning del modello LLM su un dataset specifico del dominio. Anziché usare un LLM generico, possiamo adattare un modello più piccolo (es.\ \code{Phi-3-mini}) a generare query SPARQL corrette per il Datavault specifico.

    Tecniche come LoRA (\textit{Low-Rank Adaptation}) permettono un fine-tuning efficiente senza memorizzare l'intero modello in memoria. Il modello apprende:

    \begin{enumerate}
        \item La sintassi SPARQL corretta per il Datavault;
        \item L'ontologia specifica (predicati, tipi di dato, convenzioni di naming);
        \item Strategie di generazione robuste (evitare predicati inesistenti, gestire date, etc.);
        \item Quando e come segnalare query out-of-domain.
    \end{enumerate}

    Questo approccio riduce significativamente il rischio di allucinazioni e aumenta l'accuratezza della generazione di query.

\section{Conclusioni del capitolo}\label{sec:conclusioni_cap2}
In questo capitolo abbiamo analizzato il contesto scientifico e tecnologico che motiva il nostro contributo. Abbiamo visto che:
    \begin{enumerate}
        \item Gli approcci tradizionali (SPARQL nudo, filtri rigidi) falliscono nel dominio umanistico perché impongono barriere cognitive, strutturali e pratiche;
        \item I modelli teorici dell'architettura dell'informazione (berrypicking, precision e recall, carico cognitivo) giustificano il ricorso a un'interfaccia conversazionale mediata da LLM;\@
        \item Le tecnologie emergenti (RAG, fine-tuning con LoRA) rendono possibile un'implementazione affidabile e accurata di tale interfaccia.
    \end{enumerate}

    Avendo chiarito il fondamento scientifico e tecnologico della nostra soluzione, il \cref{chap:prototipo} presenta nel concreto il nostro LLM Query Builder: il dominio di applicazione (DH.ARC e il Datavault), il profilo degli utenti target, l'architettura funzionale del sistema, e i casi d'uso che dimostrano il valore aggiunto della nostra proposta nel contesto umanistico.